\documentclass[instructions.tex]{subfiles}
\begin{document}
 
\chapter{Exemples de chasteté.}
 
La chasteté n’est pas si rare que les libertins le pensent. Sans compter l’exemple de l’incomparable Marie, mère du Sauveur, qui eût mieux aimé n’être jamais mère de Dieu que de cesser d’être vierge, combien de millions de personnes ont consacré et consacrent encore aujourd’hui à Dieu leur chasteté !
 
\primo Combien de vierges ont donné leur vie pour cette vertu ! Sainte Agathe, pour conserver son innocence, ne se laissa-t-elle pas déchirer et couper le sein ? Sainte Agnès, à l’âge de douze ans, ne souffrit-elle pas l’activité du feu, aimant mieux être brûlée vive que de perdre sa pureté ? Saint Pélage ne souffrit-il pas avec courage d’être taillé en pièces pour le même sujet ? Sainte Potamienne ne souffrit-elle pas les huiles brûlantes plutôt que de consentir à la passion de son maître ? Le jeune prince saint Casimir ne dit-il pas qu’il aimait mieux mourir que de prendre un remède qui pouvait être dangereux pour sa chasteté\footnote{Malo mori quàm fœdari. S. Cas.}.
 
On sait ce que firent les vierges de la ville de Saint-Jean-d’Acre lorsque cette ville fut prise d’assaut. Ces saintes filles se firent des plaies au visage, se coupèrent les lèvres, afin que les soldats, les voyant défigurées, en eussent horreur et n’attentassent pas à leur pureté. Rougissez, chrétiens, de votre lâcheté, en voyant ces héros et ces héroïnes de la chasteté.
 
\secundo Le paganisme même fournit des exemples mémorables de cette vertu. Chez les Romains, la fidélité conjugale était si inviolable, que Lucrèce ayant été déshonorée par surprise et par force, ne put survivre à un tel affront. Elle en fit part à son époux, toute baignée de larmes. Son mari tâchant de calmer sa douleur, sur ce que cette infidélité n’était pas volontaire, elle répondit que, quoique innocente, elle ne pouvait plus vivre après avoir été déshonorée.
 
Virginie, jeune Romaine, ayant été enlevée par l’ordre d’un sénateur, le père accourut aux cris de sa fille, et lui dit : On vous enlève, ma fille, pour vous ôter votre chasteté ; mais lequel voudriez-vous choisir, de perdre la vie, ou de perdre votre pudeur ?... J’aime mieux, répondit-elle, perdre la vie. Ensuite le père, tirant un poignard, l’enfonça dans le sein de Virginie, aimant mieux voir sa fille morte que de la voir déshonorée.
 
Le sénat de Rome envoya au roi Porsenna dix jeunes filles de qualité, qu’il avait demandées en otage. La jeune Clélie fit comprendre à ses compagnes le danger où était leur pudeur dans le palais de ce prince, et leur persuada de la suivre. Clélie, s’étant jetée toute vêtue dans le Tibre, ses compagnes la suivirent, à la nage, aimant mieux exposer leur vie que d’exposer leur vertu.
 
Une vestale, accusée d’avoir perdu sa pudeur, fit connaître son innocence, et fut renvoyée ; mais le pontife la reprit de ce qu’elle avait l’air trop enjouée, et lui dit : Si vous aviez eu plus de modestie dans vos parures et dans vos manières, on vous aurait crue chaste. La dissipation et des airs si enjoués ne conviennent point à une fille qui a soin de sa pudeur et de sa réputation.
 
Ces exemples de chasteté et de modestie ne devraient-ils pas charger de honte tant de chrétiens et de chrétiennes qui aujourd’hui ont moins de retenue que les paiens ?
 
La chasteté convient à tous, sans excepter même les personnes mariées : chacun doit la pratiquer selon son état. Sans cette vertu, on ne peut plaire à Dieu ni être sauvé. Si vous avez eu le malheur de la souiller, le remède qui vous reste, c’est de vivre en chaste pénitent.
 
\section{Résolutions}
 
\primo Je ne regarderai jamais la chasteté comme une vertu impraticable.
 
\secundo Je m’y encouragerai par l’exemple de tant de héros qui l’ont portée jusqu’à la perfection.
 
\tertio Et je me dirai souvent à moi-même : Pourquoi ne pourrais-je pas, avec le secours du Ciel, ce que ceux-ci et celles-là ont pu ? 

\prayer{Mon Dieu, mettez dans mon cœur un grand amour pour la chasteté, et faites-moi la grâce de ne plus rien me permettre qui puisse blesser cette belle vertu.}
 
\end{document}
